% ── results/qmap.tex ──────────────────────────────────────────────────────────
\subsection{Homology-Aware QMAP Benchmark}
\label{sec:res_qmap}

\begin{table}[h]
\centering
\caption{QMAP homology-aware benchmark.
  Values are mean Pearson correlation across five official splits;
  JEPA results are mean $\pm$ SD over three seeds (42, 7, 123).
  Prior baselines are from the QMAP leaderboard release.}
\label{tab:qmap}
\begin{tabular}{lccc}
\toprule
\textbf{Method} & \textbf{Full E.~coli} & \textbf{High-eff. E.~coli} & \textbf{HC50} \\
\midrule
QMAP ESM2 linear baseline   & 0.360 & 0.160 & 0.070 \\
Witten \& Witten 2019        & 0.510 & 0.220 & --- \\
Cai et al.\ 2025             & \textbf{0.520} & 0.290 & --- \\
\midrule
JEPA frozen + Ridge          & 0.496 & 0.351 & --- \\
JEPA head-only               & $0.501 \pm 0.002$ & $0.372 \pm 0.005$ & $\mathbf{0.327 \pm 0.004}$ \\
JEPA conditional head        & $0.512 \pm 0.009$ & $\mathbf{0.388 \pm 0.013}$ & $0.307 \pm 0.000$ \\
\bottomrule
\end{tabular}
\end{table}

The conditional property head surpasses all reported high-efficiency E.~coli
baselines by $\mathbf{+0.098}$ PCC ($0.388$ vs.\ $0.290$).
On full E.~coli, JEPA is competitive with Cai et al.\ 2025 ($0.512$ vs.\
$0.520$; paired split-seed mean $\Delta {=} {+}0.011$, SD $0.021$).
For HC50, the task-specific head ($0.327$) outperforms the shared conditional
head ($0.307$; mean paired $\Delta {=} {-}0.021$), indicating that bacterial
MIC and mammalian haemolysis should not be treated as a shared prediction task.

\finding{Under homology-aware evaluation, JEPA-AMP is competitive on full E.~coli MIC
and distinctly stronger on high-efficiency E.~coli. The HC50 result confirms
that toxicity endpoints benefit from dedicated heads rather than shared
bacterial MIC representations.}
